\section{问题重述与统一符号}

机器人由两类定位系统给出异频观测。前三问依次要求在无噪声、含噪且含固定系统偏差、实际复杂误差三种条件下恢复统一 10 Hz 轨迹；第四问在固定轨迹上安排射击和拍照任务。

为了避免不同代码模块对“时间偏差”的正负号定义不一致，本文只在报告层使用一个符号：
\begin{equation}
  t_{2,\mathrm{aligned}}=t_2+\delta.
\end{equation}
因此 $\delta>0$ 表示方式 2 时间戳整体增加，$\delta<0$ 表示整体减小。底层程序仍保留各自已经验证的内部变量定义，自动报告脚本负责转换到统一口径。

\input{generated/table_time_offsets.tex}

四问的总体技术路线为
\[
\text{时间同步}
\longrightarrow
\text{时空标定}
\longrightarrow
\text{偏差检验与鲁棒融合}
\longrightarrow
\text{固定轨迹任务调度}.
\]

所有正式结果均写入 \texttt{05\_results/}，论文图来自 \texttt{06\_figures/}，因此论文正文不再作为数值的唯一来源。
